\inicializarSeccion{Parametrización de posición de cargas}

\tab Una vez los parámetros establecidos podemos determinar las posiciones de cada una de las cargas positivas que utilicemos, así que podemos utilizar la siguiente ecuación para describir su posición.

\begin{equation}
    (x_{pos}, y_{pos}) = \left(-\frac{L_p}{2} + (i-1) \cdot \frac{L_p}{N_p-1}, \frac{d}{2}\right)
\end{equation}

\tab De manera similar, para las cargas negativas, su posición puede ser descrita por la siguiente ecuación:

\begin{equation}
    (x_{neg}, y_{neg}) = \left(-\frac{L_n}{2} + (j-1) \cdot \frac{L_n}{N_n-1}, -\frac{d}{2}\right)
\end{equation}

\tab Gracias a este modelo podemos trabajar con múltiples cargas distribuidas a lo largo de nuestro campo eléctrico.